% Preamble for a minimal theoretical CS / graph theory style
\usepackage[T1]{fontenc}
\usepackage[utf8]{inputenc}
\usepackage{lmodern}
\usepackage{microtype}

\usepackage{amsmath,amssymb,amsthm,mathtools}
\usepackage{geometry}
\usepackage{hyperref}

\geometry{margin=1in}

% \setlength{\parindent}{0pt}
% \setlength{\parskip}{0.4em}

\hypersetup{
  colorlinks=true,
  linkcolor=black,
  citecolor=black,
  urlcolor=black
}




\usepackage{physics}
\usepackage[shortlabels]{enumitem}
\usepackage{wrapfig, float}
\usepackage{tikz}
\usetikzlibrary{calc}
\usetikzlibrary{positioning, shapes.geometric, fit}




\DeclarePairedDelimiter\floor{\lfloor}{\rfloor}
\usepackage{thmtools}



\usepackage[justification=centering, margin=30pt, labelfont=bf]{caption}
% \usepackage{mathabx}
\usetikzlibrary{patterns}
\usetikzlibrary{arrows}



\renewcommand{\and}{\hspace{6pt}\text{and}\hspace{6pt}}


\newcommand*{\heading}[1]{\begin{center}{\huge\textbf{#1}}\end{center}}
\newcommand*{\subheading}[1]{\begin{center}{\large{#1}}\end{center}}


\newcommand{\Or}{\hspace{6pt}\text{or}\hspace{6pt}}


\DeclareMathOperator{\Log}{Log}
\DeclareMathOperator{\sg}{sg}
\renewcommand{\rm}{\text{rm}}
\DeclareMathOperator{\qt}{qt}
\renewcommand{\Pr}{\text{Pr}}
\DeclareMathOperator{\id}{id}
\DeclareMathOperator{\Sub}{Sub}
\DeclareMathOperator{\cut}{cut}
\DeclareMathOperator{\Fr}{Fr}
\DeclareMathOperator{\Bin}{Bin}
\DeclareMathOperator{\supp}{supp}
\DeclareMathOperator{\DP}{DP}
\DeclareMathOperator{\LFP}{LFP}
\DeclareMathOperator{\up}{\uparrow}
\DeclareMathOperator{\down}{\downarrow}
\DeclareMathOperator{\CSP}{CSP}
\DeclareMathOperator{\diam}{diam}
\DeclareMathOperator{\Cl}{Cl}
\DeclareMathOperator{\fl}{fl}
\DeclareMathOperator{\mon}{mon}
\DeclareMathOperator{\poly}{poly}
\DeclareMathOperator{\Hom}{Hom}
\DeclareMathOperator{\Mod}{Mod}
\DeclareMathOperator{\CFI}{CFI}
\DeclareMathOperator{\Core}{Core}
\DeclareMathOperator{\td}{td}
\DeclareMathOperator{\tw}{tw}
\DeclareMathOperator{\Undef}{undef}


\newcommand{\If}{\hspace{6pt}\text{if}\hspace{6pt}}

\newcommand{\A}{\mathcal{A}}
\newcommand{\C}{\mathcal{C}}
\newcommand{\D}{\mathcal{D}}
\newcommand{\F}{\mathcal{F}}
\newcommand{\K}{\mathcal{K}}
\newcommand{\T}{\mathcal{T}}
\renewcommand{\H}{\mathcal{H}}
\newcommand{\M}{\mathcal{M}}
\newcommand{\Z}{\mathbb{Z}}
\newcommand{\imm}{\preccurlyeq_{\text{imm}}}
\newcommand{\immsimp}{\preccurlyeq_{\text{immsimp}}}
\newcommand{\revimm}{\succcurlyeq_{\text{imm}}}
\newcommand{\immo}{\preccurlyeq_{\text{imm-odd}}}
\newcommand{\oddo}{\xrightarrow{\text{oddo}}}




\newcommand{\B}{\mathbb{B}}
\newcommand{\N}{\mathbb{N}}
\renewcommand{\flat}{\text{flat}}


\newcommand\restr[2]{{% we make the whole thing an ordinary symbol
  \left.\kern-\nulldelimiterspace % automatically resize the bar with \right
  #1 % the function
  \vphantom{\big|} % pretend it's a little taller at normal size
  \right|_{#2} % this is the delimiter
  }}




% \setlength{\parindent}{0pt}

\newtheorem{theorem}{Theorem}
\newtheorem{lemma}[theorem]{Lemma}
\newtheorem{corollary}[theorem]{Corollary}
\newtheorem{proposition}[theorem]{Proposition}
\newtheorem*{note}{Note}
\newtheorem{question}[theorem]{Question}
\newtheorem{conjecture}[theorem]{Conjecture}

\theoremstyle{definition}
\newtheorem{definition}[theorem]{Definition}
\newtheorem{example}[theorem]{Example}



% \newcounter{claimcount}
% \newenvironment{claim}{
%   \refstepcounter{claimcount}
%   \begin{trivlist}
%   \item[\hskip \labelsep {\textit{Claim \theclaimcount.}}] \itshape
% }{
%   \end{trivlist}
% }
% \newenvironment{proofofclaim}[1][Proof of claim]{
%   \begin{proof}[#1]
%   \renewcommand{\qedsymbol}{$\lozenge$} % Changes the QED symbol to a hollow diamond
% }{
%   \end{proof}
% }


\newenvironment{claim}{
  \refstepcounter{theorem} % Increments the theorem counter instead
  \begin{trivlist}
  \item[\hskip \labelsep {\textit{Claim \thetheorem.}}] \itshape
}{
  \end{trivlist}
}

\newenvironment{proofofclaim}[1][Proof of claim]{
  \begin{proof}[#1]
  \renewcommand{\qedsymbol}{$\lozenge$} 
}{
  \end{proof}
}



